\documentclass[12pt,a4paper]{article}

\usepackage{cmap}
\usepackage[T2A]{fontenc}
\usepackage[utf8]{inputenc}
\usepackage[russian]{babel}
\usepackage{amsmath,amssymb}
\usepackage[a4paper,margin=2cm]{geometry}

\setlength{\parindent}{0pt}
\setlength{\parskip}{0.45em}
\sloppy

\begin{document}

\begin{center}
{\Large\bfseries Билет 8: вопросы для проверки знаний}
\end{center}

\noindent\fbox{%
  \begin{minipage}{\dimexpr\linewidth-2\fboxsep-2\fboxrule\relax}
  \normalfont
Геометрические приложения определенного интеграла: площадь криволинейной
трапеции, объем тела вращения, длина кривой. Площадь поверхности вращения (без
доказательства).
  \end{minipage}%
}

\section*{Вопросы на понимание}

\begin{enumerate}
\item Почему в формуле площади криволинейной трапеции требуется рассматривать неотрицательную функцию \(f\)?
\item Чем площадь криволинейной трапеции под графиком неотрицательной функции отличается от алгебраического значения определенного интеграла для функции, меняющей знак?
\item Как нижние и верхние клеточные приближения трапеции связаны с нижними и верхними суммами Дарбу?
\item Почему интегрируемость функции \(f\) позволяет перейти от клеточных приближений к точной площади трапеции?
\item Почему в формуле объема тела вращения появляется \(\pi f^2(x)\), а не просто \(f(x)\)?
\item Какой геометрический смысл имеет сечение тела вращения плоскостью \(x=\mathrm{const}\)?
\item Проверьте формулы площади трапеции, объема тела вращения, длины графика и площади поверхности вращения на примере постоянной функции \(f(x)=R>0\).
\item Почему длину кривой определяют через супремум длин вписанных ломаных, а не через расстояние между ее концами?
\item Почему одной непрерывности параметризации недостаточно, чтобы автоматически пользоваться формулой длины \(\int_a^b |r'(t)|\,dt\)?
\item Откуда в формуле длины графика \(y=f(x)\) берется множитель \(\sqrt{1+(f'(x))^2}\)?
\item Чем концептуально отличаются формулы для объема тела вращения и площади поверхности вращения?
\item Зачем в определении площади поверхности вращения сначала вращают вписанные ломаные и затем переходят к пределу?
\end{enumerate}

\section*{Источники для сверки}

\texttt{target/answer\_question\_08.tex};
\texttt{IvGE\_1.pdf}, стр. 146--150, 230--238;
дополнительно просмотрено: \texttt{IvGE\_2.pdf}, стр. 171--173.

\end{document}
